%!TEX root = ../../main.tex
\section{평가기 의존성과 점수 분해의 읽기}
\label{sec:disc-evaluator}

$\Ysvi$가 $\dV$의 부호이므로 모든 예측 가능성 주장은 동결 fit85 평가기와 결과 시점에 조건부이다. 킬과 오브젝트 결과와의 대응은 $\SVI$가 자의적 잡음이 아님을 지지하지만, 입력의 공유와 공통 원인 때문에 그 표를 인과적 교전 효과나 $q$의 정확도로 읽을 수는 없다. 방향 AUC를 높이기 위해 $V$를 재조정하지 않은 것은 의도적이다. 그렇게 하면 추정 대상 자체가 바뀐다. 동료 평가기와의 약 9\% 부호 불일치는 측정의 한계로 남으며, 이 논문의 버전 경계 안에서 아키텍처를 계속 찾을 근거가 아니다.

15.16에서 $q$와 $\PTflex$의 Brier 격차는 $q$의 더 큰 판별 성분과 다소 더 큰 오보정 성분과 함께 움직이고, KR/NA1 16.13에서는 판별이 여전히 $q$에 유리하면서 오보정이 순 손실을 지배한다. 이는 평가 표본에서 적정 점수가 어떻게 분해되는지의 기술이다. 게임 클라이언트 안의 인과 기전을 식별하지 않으며, 저차원 재보정기가 전이를 복구한다는 것을 증명하지도 않는다. 어댑터 실험은 동결 규칙 아래 별도의 v2 계열로 이연된다(부록 \ref{app:deferred} F1).
